\begin{figure}[!ht]
\centering
\begin{tikzpicture}[x=0.028cm, y=0.030cm, >=Latex,
  lab/.style={font=\scriptsize, align=left},
  lead/.style={black!55, thin, shorten >=2pt}]
  % axes
  \draw[->, thick] (-8,0) -- (262,0) node[below] {\footnotesize $a_{R_1}$};
  \draw[->, thick] (0,-8) -- (0,208) node[left] {\footnotesize $a_{R_2}$};
  \foreach \x in {50,100,150,200} \draw (\x,0) -- (\x,-5) node[below] {\tiny \x};
  \foreach \y in {50,100,150} \draw (0,\y) -- (-5,\y) node[left] {\tiny \y};

  % loads produced by the demand-derived price sweep (the seed pool)
  \foreach \p in {(0,127),(0,145),(0,163),(0,168),(0,182),(7,168),(14,168),
                  (16,100),(19,117),(21,168),(22,134),(25,151),(28,126),
                  (28,168),(142,0),(166,0),(190,0),(214,0),(238,0)}
      \fill[black!55] \p circle[radius=2.6];

  % upper-right frontier of the convex hull
  \draw[black!75, dashed, thick] (0,182) -- (28,168) -- (238,0);
  \node[lab, black!75, fill=white, inner sep=1.5pt, rotate=-31] at (168,64) {frontier of $\operatorname{conv}(\Omega^{s})$};

  % dominating blend, and the load an integer plan needs
  \fill[blue!70!black] (74,131.2) circle[radius=3];
  \node[star, star points=5, star point ratio=2.2, fill=red!75!black,
        inner sep=1.7pt] at (74,126) {};
  \draw[blue!70!black, <->, thin] (74,129) -- (74,128.2);

  \node[lab, blue!70!black, anchor=west, text width=40mm] (bl) at (86,198)
      {$0.78\,(28,168)+0.22\,(238,0)$\\ $=(74,131.2)$: the blend the\\
       relaxation buys};
  \draw[lead, blue!70!black] (bl.south west) -- (78,133);

  \node[lab, red!75!black, fill=white, inner sep=1.5pt, anchor=west, text width=38mm] (tg) at (96,104)
      {$(74,126)$: the load a\\ twelve-truck plan needs};
  \draw[lead, red!75!black] (tg.west) -- (79,124);

  % a dual vector, drawn inside the hull where there is room
  \draw[->, thick] (108,52) -- ++(24,18);
  \node[lab, anchor=north east] at (108,52) {$\pi$};
  \node[lab, anchor=west, text width=36mm] at (56,24)
      {for \emph{every} $\pi$, the maximum of $\pi^{\!\top}\!a$ is attained
       at a vertex of the frontier};
  % inset: the blend and the load an integer plan needs are 5.2 units apart on a
  % 208-unit axis, so they read as one mark at print size.
  \begin{scope}[shift={(176,138)}, x=0.30cm, y=0.32cm]
     \draw[black!35] (-1.2,-1.2) rectangle (6.2,4.2);
     \draw[black!75, dashed] (-1.2,3.4) -- (6.2,2.2);
     \fill[blue!70!black] (2.5,2.9) circle[radius=0.16];
     \node[star, star points=5, star point ratio=2.2, fill=red!75!black,
           inner sep=1.1pt] at (2.5,0.9) {};
     \node[lab, anchor=west, scale=0.9] at (3.1,2.9) {$(74,131.2)$};
     \node[lab, anchor=west, scale=0.9] at (3.1,0.9) {$(74,126)$};
     \node[lab, anchor=south, scale=0.9] at (2.5,4.2) {inset, $\times 20$};
  \end{scope}
\end{tikzpicture}
\caption{Why the load a cheaper plan needs is never priced. Dots are
the supplier-bound loads obtained from a sweep of demand-derived price
vectors; the dashed line is the upper-right frontier of their convex
hull. The pricing objective is linear, so for \emph{any} dual vector
$\pi$ its optimum is attained at a vertex of that frontier. The load
$(74,126)$ lies strictly below it --- it is dominated componentwise by
the blend $0.78\,(28,168)+0.22\,(238,0)=(74,131.2)$ --- so no choice of
duals returns it. The relaxation buys that blend for one truck, while
an integer plan must pay for both of the trucks it mixes. The construction of
Section~\ref{sec:relaxfix-flow} reaches the interior directly: it
partitions a delivery vector into legal loads instead of asking a
linear objective for them.}
\label{fig:hull}
\end{figure}
